\begin{table}[tbp] \centering
\newcolumntype{R}{>{\raggedleft\arraybackslash}X}
\newcolumntype{L}{>{\raggedright\arraybackslash}X}
\newcolumntype{C}{>{\centering\arraybackslash}X}

\caption{Leave-one-out: the wine result survives dropping any single canton (col 5)}
\label{tab:T_loo_gate}
\begin{tabularx}{\linewidth}{@{}lCC@{}}

\toprule
{ }&{FL AME (pp)}&{OLS coef (HC3)} \tabularnewline
\midrule \addlinespace[\belowrulesep]
Full sample (N=25)&0.437 (0.194)&0.473 (0.307) \tabularnewline
Drop Neuchatel (NE)&0.443 (0.215)&0.482 (0.596) \tabularnewline
Drop Vaud (VD)&0.497 (0.273)&0.532 (0.536) \tabularnewline
LOO range [min, max]&[0.321, 0.653]&[0.353, 0.697] \tabularnewline
\bottomrule \addlinespace[\belowrulesep]

\end{tabularx}
\\ \parbox{\linewidth}{\footnotesize Yes-vote share (vote \#68) on national wine revenue share, col-5 spec (+ French share + Absinthe trade share + Protestant share + log density). FL AME = fractional-logit average marginal effect, percentage points of yes-share per percentage point of wine share; OLS coef = HC3, on the 0-100 yes-share scale. Standard errors in parentheses. Each row re-estimates dropping the named canton(s). LOO range spans all 25 single-canton deletions (FL AME min at drop-ZH, max at drop-FR). The wine effect stays positive and similar in magnitude under every single-canton deletion, including drop-Vaud: it is not driven by any one canton.}
\end{table}
